% Folder 79, batch 06 — archivists' pages 101–120.
% Pass: Opus 5.5 (claude-opus-5-5), 2026-10-03 — first pass, unchecked against the pages by a human.
%
% Content: end of a paragraph on the torsor \vec\pi under \Lambda (formulas 71–74), then
% Grothendieck's section VII « Structures triangulaires congruentielles » (his pp. 52–61):
% triangular structures (S,F), axioms T1–T3 and T_0 1–T_0 6, the centraliser Z' of G_0^+,
% the graphs \Gamma_\lambda and polyhedra P_\lambda. No page skipped.
\documentclass[11pt,a4paper]{article}
\input{../preamble/grothendieck}

\folder{79}
\batch{6}
\pages{101}{120}
\foldertitle{2-polyèdres réguliers triangulés, et modules libres de rangs 2 sur les anneaux $\Lambda=\mathbb{Z}/n\mathbb{Z}$ et $\Lambda=\mathbb{Z}$ : notes manuscrites (s.d.).}
\dating{[à partir de 1980]}
\watermark{Édition de démonstration}

\begin{document}

\note{Les pp.~101 sqq.\ reprennent un argument commencé avant ce lot : les formules (71) à (74) de la p.~101 achèvent un paragraphe antérieur.}

\page{101}
\note{En tête, le numéro de l'auteur, partiellement couvert d'encre : « 5\ill{}* ».}

$\exists !\, \lambda_{e'} \in \Lambda$ tel qu'on ait
\[
(71) \qquad \lambda_{\Gamma,s} = \lambda_{e'} \wedge_{\mu_2} (e \wedge e')
\]
et on a bien sûr
\[
(72) \qquad \lambda_{\bar e'} = -\lambda_{e'} \quad \text{si } \mu(i') = \{e', \bar e'\}.
\]
Ainsi, $\{\lambda_{e'}, \lambda_{\bar e'}\}$ est l'image inverse de $\lambda \in \Lambda^{*}/\pm 1$ dans $\Lambda^{*}$, et il existe un unique $e' \in \mu(i')$, tel qu'on ait
\[
(73) \qquad \lambda = \lambda_{e'}, \quad \text{i.e.} \quad \lambda_{\Gamma,s} = \lambda \wedge (e \wedge e') \in \Gamma_{i,i'} .
\]
A priori, ce choix de $e'$ dépend de celui de $s \in S$ -- en fait il est indépendant de ce choix. On trouve ainsi l'opération \uncertain{sur} \struck{\ill{}} $\vec\pi$, qui fait de $\vec\pi$ un torseur sous $\Lambda$.

L'application
\[
(74) \qquad \tau_F : \underbrace{\vec F}_{\vec F_M} \longrightarrow \vec\pi
\]
qui relève $\tau : F \to \pi$ est maintenant \uncertain{évidente}, \add{orientation}\note{mot ajouté sous la ligne, sans place nette dans la phrase.} \struck{\ill{}} \uncertain{donnée} \uncertain{par} $f \in \vec F$ \uncertain{définit} \uncertain{à la fois} un $\tau(f) = i \in \pi$, et un $e \in \mu(i) = \operatorname{Or} P_i$ \ill{} une orientation de la structure triangulaire ; lesquelles \ill{} \ill{} \ill{}.

\section{VII Structures triangulaires congruentielles}
\note{Le titre est celui de l'auteur, en tête de la p.~102, souligné ; il ouvre un nouveau paragraphe numéroté VII, dont la numérotation des formules repart à (1).}

\page{102}
Une \emph{structure triangulaire} est pour nous un couple
\[
(1) \qquad (S, F) \qquad
\begin{array}{l}
S \text{ un ensemble (« \emph{sommets} »),} \\
F \subset \mathfrak{P}_3(S) \ \text{(l'ens.\ des \emph{faces})}
\end{array}
\]
satisfaisant à l'axiome
\[
(2) \quad \text{(T1)} \qquad S = \bigcup_{f \in F} f \qquad \text{donc} \quad S \neq \emptyset \iff F \neq \emptyset .
\]
\note{Le repère « (T1) » est d'abord écrit puis biffé à gauche, et récrit, cerclé, sous la ligne.}
Les structures qui nous intéressent satisfont à des axiomes supplémentaires, inspirés par le cas des $(S_M, F_M)$ où $M$ est un module libre de rang~2 sur un $\mathbb{Z}/n\mathbb{Z}$, $n \geq 3$. Posons d'ailleurs
\[
(3) \qquad A = \{ a \in \mathfrak{P}_2(S) \mid \exists f \in F,\ f \supset a \}
\]
(ens.\ des \emph{arêtes}) ; on demande
\[
(4) \quad \text{(T2)} \qquad \forall a \in A, \text{ il existe exactement deux } f \in F \text{ telles que } f \supset a
\]
(toute arête \uncertain{incidente} à exactement deux faces).

En d'autres termes, le triple
\[
(5) \qquad \Gamma = (F, A, R_2 \subset F \times S), \qquad R_2 = \{ (f, a) \in F \times A \mid f \supset a \}
\]
\note{Le $R_2$ et les indices de $F$ et $A$ sont surchargés ; « $\subset F \times S$ » est lu tel qu'écrit, bien que $R_2$ soit défini dans $F \times A$.}
est une structure de \emph{graphe strict}, ayant \ill{}

\page{103}
\note{Pagination de l'auteur : 52.}
$F$ comme ens.\ des sommets, et $A$ comme ens.\ d'arêtes [\struck{\ill{}} \struck{graphe strict} : toute arête est \uncertain{une} \uncertain{famille} et déterminée par l'ens.\ de ses extrémités] -- ici les sommets du graphe sont d'ordre~3 (\struck{\ill{}} \add{les arêtes issues} d'un « sommet » $f$ du graphe correspondant aux trois parties à deux éléments de $f \in \mathfrak{P}_3(S)$).

On pose
\[
(6) \qquad \pi = \pi_0(\Gamma) \qquad \text{ens.\ des composantes connexes du graphe } \Gamma
\]
\struck{d'où \ill{} \ill{} $i \in \pi$} d'où des \uncertain{projections}, surjections
\[
(7) \qquad
\left\{
\begin{array}{l}
F \xrightarrow{\ \tau_F\ } \pi \\[2pt]
A \xrightarrow{\ \tau_A\ } \pi
\end{array}
\right.
\]
dont les fibres en $i \in \pi$ sont notées $F_i$, $A_i$, donc
\[
(8) \qquad
\left\{
\begin{array}{ll}
F = \coprod_{i \in \pi} F_i & \text{les } F_i \neq \emptyset \\[2pt]
A = \coprod_{i \in \pi} A_i & \text{les } A_i \neq \emptyset
\end{array}
\right.
\]
et

(9) si $a \in A$, $f \in F$, $a \subset f$, alors $f \in F_i \iff a \in A_i$ ; si $f, f' \in F_i$, il existe une chaîne $f_0 = f, f_1, \ldots, f_N = f'$ de faces \struck{\ill{}} \add{\uncertain{deux à deux} consécutives} adjacentes (i.e.\ $f_i \cap f_{i+1} \in A$ i.e.\ $\operatorname{card}(f_i \cap f_{i+1}) = 2$) qui les joignent.

\page{104}
Notons qu'on a sur le triple
\[
(10) \qquad
\left\{
\begin{array}{lll}
P = (S, A, F) & & \text{des relations « d'incidence »} \\[2pt]
R_1 \subset S \times A & & R_1 = \{ (s, a) \in S \times A \mid s \in a \} \\[2pt]
R_2 \subset A \times F & & R_2 = \{ (a, f) \in A \times F \mid a \subset f \} \\[2pt]
R_3 = R_2 \circ R_1 \subset S \times F & & R_3 = \{ (s, f) \in S \times F \mid s \in f \} ,
\end{array}
\right.
\]
qui satisfont « l'axiome canonique » des structures 2-polyédrales
\[
(11) \qquad
\left\{
\begin{array}{l}
\text{a) } \forall a \in A, \text{ il existe exactement deux sommets incidents} \\
\phantom{\text{a) }} \text{et deux faces incidentes} \\[2pt]
\text{b) } \forall (s, f) \in R_3, \text{ il existe exactement deux arêtes} \\
\phantom{\text{b) }} \text{incidentes à la fois à } s \text{ et à } f
\end{array}
\right.
\]
et aussi l'axiome de connexité sur les faces -- qui sont en effet des triangles.
\marginal{\emph{NB} $R_2'$ de (5) est la relation \uncertain{symétrique} de $R_2$ de (10).}

Considérons d'autre part la question de connexité en un sommet $s \in S$, \uncertain{ou} son « étoile » $\operatorname{Et}(s)$ ou $A(s)$. Introduisons
\[
(12) \qquad A(s) = \{ t \in S \mid \{s, t\} \in A \} ,
\]
\note{Sous (12), une annotation dense, en partie biffée et accolée : « l'un des sommets \ill{} d'une arête \struck{\ill{}} \uncertain{contenant} $s$ », puis, sous une accolade, « \uncertain{pour décrire comme} arêtes \ill{} ». Le symbole de gauche, d'abord écrit $\operatorname{Et}(s)$, est surchargé en $A(s)$.}
$P(s)$, \ill{}
\[
F(s) = \bigl\{ \{t, t'\} \in \mathfrak{P}_2(S) \bigm| \{s, t, t'\} \in F \bigr\} .
\]
\note{La ligne « $P(s)$, \ill{} » est en partie écrite sous une accolade ; la notation de la face $F(s)$ est surchargée sur un symbole biffé.}
\marginal{$P(s)$ peut être considéré comme \ill{} de la structure triangulaire $P$ en $s$ \uncertain{(13)}.}
\[
(13) \qquad R_3(s) = \{ f \in F \mid s \in f \}
\]
\note{Numéro « 21 » écrit au-dessus de $R_3(s)$, sans rapport visible avec la numérotation des formules.}
On a bien sûr des décompositions

\page{105}
\note{Pagination de l'auteur : 53, cerclé.}
\[
(14) \qquad
\left\{
\begin{array}{ll}
A(s) = \coprod_{i \in \pi} A_i(s) & A_i(s) = \{ t \in S \mid \{s, t\} \in A_i \} \\[4pt]
F(s) = \coprod_{i \in \pi} F_i(s) & F_i(s) = \bigl\{ \{t, t'\} \in \mathfrak{P}_2(S) \bigm| \{s, t, t'\} \in F_i \bigr\}
\end{array}
\right.
\]
et correspondante : une décomposition des contours $P(s)$ en
\[
(15) \qquad P(s) = \coprod_{i \in \pi} P_i(s) .
\]
Pour l'axiome de connexité en $s$ on \uncertain{vérifie} que si tous les $A_i(s)$ sauf un seul sont vides ! \struck{Ceci} En fait, on ne suppose pas cet axiome satisfait, mais on introduit les
\[
(16) \qquad P_i = (S, A_i, F_i, R_{1i}, R_{2i}, R_{3i} = R_{2i} \circ R_{1i})
\]
qu'on interprète comme des structures triangulaires sur le même ens.\ de sommets $S$, et on exige l'axiome ci-dessous
\[
(17) \quad \text{(T3)} \qquad \forall i \in \pi, \ (S, F_i) \text{ est une structure triangulaire,}
\]
i.e.\ $\bigcup_{f \in F_i} f = S$ (satisfaisant alors évidemment l'axiome \uncertain{conn.}\ et (T2)) et \emph{stricte} au sens qu'en plus de (T1) (T2), on a connexité des contours $P_i(s)$.

Cela s'exprime donc aussi par les deux conditions

\page{106}
(18)
\begin{itemize}
\item T3 a) $\forall f \in F$ et $s \in S \setminus f$, il existe une « \struck{\ill{}} \add{chaîne} » de faces $f_0 = f, f_1, \ldots, f_N$ (deux à deux consécutives bien sûr adjacentes) avec $s \in f_N$ ;
\item T3 b) Si $f, f' \in F$ \struck{$\{s,t\}, \{s,t'\}$} sont deux \struck{arêtes adjacentes} \add{faces contenant} un \struck{\ill{}} sommet $s$, qui peuvent être jointes par une « chaîne » de faces $f_0 = f, \ldots, f_N = f'$, alors on peut trouver une telle chaîne formée de faces contenant $s$.
\end{itemize}
\note{Dans (18) b), le $f_0$ est surchargé ; le renvoi « (18) » est écrit en marge à gauche.}

Une structure $(S, F)$ satisfaisant T1 à T3 peut aussi s'interpréter comme une structure
\[
(19) \qquad \bigl( S, (F_i)_{i \in \pi} \bigr) \quad \text{où} \quad
\begin{array}{l}
\pi \subset \mathfrak{P}(\mathfrak{P}_3(S)) \\
\text{est un ensemble} \\
\text{formé de parties} \\
F_i \subset \mathfrak{P}_3(S)
\end{array}
\]
\note{Sous $(F_i)_{i \in \pi}$, une accolade porte : « notation \uncertain{indicielle} pour un ensemble $\pi$ d'objets $F_i$ ».}
avec le seul axiome que chaque $F_i$ définit sur $S$ une structure de 2-polyèdre triangulé \uncertain{connexe} [tout $a \in \mathfrak{P}_2(S)$ contenu dans un $f \in F_i$ est contenu dans exactement deux tels $f$, plus axiome de connexité dans tout \struck{point} \add{sommet} $s \in S$, plus axiome de connexité globale…].

\marginal{\emph{NB} $\pi \neq \emptyset \iff S \neq \emptyset$. Par la suite, on supposera $S \neq \emptyset$, \uncertain{i.e.} (le cas où $\operatorname{card} S = 4$ \ill{} que $\operatorname{card} S \geq 4$ \uncertain{est celui du tétraèdre}, $\operatorname{card} \pi = 1$, $f = \mathfrak{P}_3(S)$…).}
\note{Note marginale écrite en oblique dans le coin inférieur gauche ; lecture très incertaine dans sa partie médiane.}

Les axiomes qui suivent sont distincts, entre autres, à assurer que \struck{\ill{}} l'ens.\ $\pi$ des structures polyédrales sur $S$ est déterminé uniquement par la connaissance de l'\emph{une} des

\page{107}
\note{Pagination de l'auteur : 54.}
structures. On va donc \struck{fixer une} \ill{} l'ens.\ $\pi$ des structures 2-polyédrales sur $S$, et \struck{retrouver} partir d'une situation où on a \emph{une} telle structure
\[
(20) \qquad (S, F_0), \qquad
\begin{array}{l}
S \text{ ensemble non vide} \\
F_0 \subset \mathfrak{P}_3(S) \\
\text{faisant de } S \text{ un 2-polyèdre triangulé,}
\end{array}
\]
par la suite, reconstruire : $(S, F)$, \uncertain{en} supposant par ailleurs particulier $F_0 = F_{i_0}$, où $i_0 \in \pi$ est un élément particulier de $\pi = \pi_0(\Gamma)$. \struck{\ill{}} Les axiomes que je vais \uncertain{exiger} pour \uncertain{telles} structures seront \struck{\ill{}} \add{distingués} \uncertain{grâce} \uncertain{par} \uncertain{l'indice}~$0$.
\[
(21) \quad (\text{T}_0 1) \qquad
\begin{array}{l}
\text{Le 2-polyèdre } (S, F_0) \text{ est \emph{régulier},} \\
\text{i.e.\ le groupe } G_0 \text{ de ses automorphismes} \\
\text{est transitif. Cela signifie} \\
\text{transitif sur l'ens.\ des repères } \operatorname{Rep}(P_0)
\end{array}
\]
où on désigne par $P_0$ la structure polyédrale
\[
(22) \qquad P_0 = (S, A_0, F_0, R_{1,0}, R_{2,0}, R_{3,0})
\]
déduite de (20), et où on a \uncertain{écrit} $\operatorname{Rep}(P_0)$ comme
\[
(23) \qquad \operatorname{Rep}(P_0) = \{ (s, t, u) \in S^3 \mid \{s, t, u\} \in F_0 \} .
\]
\note{Devant « $\text{T}_0 1$ », comme ensuite devant $\text{T}_0 2$ et $\text{T}_0 3$, un premier repère est biffé. Les numéros des deux dernières formules sont surchargés et se lisent mal (« (27) », « (33) » ?) ; la suite (20), (21) impose (22) et (23).}

\page{108}
\[
(24) \quad (\text{T}_0 2) \qquad P_0 \text{ est orientable}
\]
On désigne par $\mu_0$ l'ensemble de ses deux orientations
\[
(25) \qquad
\left\{
\begin{array}{l}
\mu_0 = \operatorname{Or}(P_0) , \\[2pt]
G_0^{+} = \operatorname{Ker}(G_0 \longrightarrow \mathfrak{S}_{\mu_0}) \\
\text{groupe des automorphismes \emph{directs} du polyèdre } P_0
\end{array}
\right.
\]
ainsi $P_0$ est simplement transitif sur l'ens.\ $\vec A_0$ des arêtes orientées de $P_0$.

On introduit
\[
(26) \qquad Z' \subset \mathfrak{S}_S, \qquad Z' = \text{sous-groupe de } \mathfrak{S}_S \text{ commutant à } G_0^{+} .
\]
$S$ est un espace homogène sous $G_0^{+}$ ; choisissons une « origine » $s \in S$ dans cet espace homogène, et introduisons le stabilisateur $G_{0s}^{+}$ de $s$ dans $G_0^{+}$, on a
\[
(27) \qquad Z' \simeq \operatorname{Norm}_{G_0^{+}}(G_{0s}^{+}) / G_{0s}^{+} .
\]
\struck{Remarque} Notons \add{(si $n$ est l'ordre des sommets de $P_0$,)} que le groupe $G_{0s}^{+}$ est isomorphe à $\mathbb{Z}/n\mathbb{Z} = \Lambda_n$ \struck{\ill{}} -- l'isomorphisme
\note{La numérotation (26) et (27) est restituée d'après le renvoi « (26) » de la note marginale et le « (27) » en face de la formule ; le corps de (26) n'a pas de numéro lisible.}

\marginal{Je \uncertain{signale} \uncertain{aussi} que \uncertain{(24)}, si \uncertain{elle} \ill{} \ill{} condition \ill{} \ill{} l'\ill{} \ill{} \ill{} \ill{} \ill{}. D'abord \uncertain{pour} \ill{} \ill{} \ill{} \ill{} orientations \ill{} $G_0$ \ill{} \ill{} (25) \ill{} $G_0^{+}$ \ill{}.}
\marginal{\uncertain{Il} \uncertain{est} \uncertain{important} \uncertain{que} l'axiome $\text{T}_0 2$ \uncertain{pour} (26) \uncertain{plus} \ill{} \ill{} \ill{} \ill{} \ill{} \ill{} \ill{} $Z'$ \ill{} \ill{} \ill{} \ill{} \ill{} $G^{0}$ \ill{} $G_0^{+}$ \ill{} \ill{} \ill{} \ill{}.}
\note{Deux longues notes marginales écrites en oblique dans l'angle supérieur gauche et le long du bord gauche ; seules des bribes se lisent.}

\page{109}
\note{Pagination de l'auteur : 55.}
dépendant du choix d'une orientation $e$ de $P_0$ -- de façon intrinsèque on a
\[
(28) \qquad G_{0s}^{+} \simeq \Lambda \wedge_{\mu_2} \mu_0
\]
\marginal{le sens des deux symboles $\mu_2$ et $\mu_0$ est tout à fait différent !}
mais l'opération naturelle de $Z'$ sur $G_{0s}^{+}$, \struck{\ill{}} $\Lambda_n$ -- module libre de rang un, fait voir une « \uncertain{opération} » canonique
\[
(29) \qquad Z' \xrightarrow{\ \chi_s\ } \Lambda_n^{*} ,
\]
\uncertain{On} peut songer à utiliser $\chi$ pour identifier $Z'$ à $\Lambda_n^{*}$ via un sous-groupe de $\Lambda_n^{*}$. Mais \uncertain{doit-il} \ill{} \ill{} \ill{}… nous \uncertain{conduit} à
\[
Z' \simeq \Lambda_n^{*} / \pm 1
\]
et $\chi_s$ sera défini par la formule
\[
(30) \qquad \chi_s(\dot\lambda) = \lambda^2 \qquad \text{pour } \lambda \in \Lambda_n^{*}
\]
\marginal{formule (*) \uncertain{v.} \uncertain{p.~46'}}
qui montre qu'on ne doit pas s'attendre que (29) soit injectif -- donc il faut trouver un autre principe pour construire un iso.\ $Z' \simeq \Lambda_n^{*} / \pm 1$.

\struck{On suppose maintenant}
\[
\text{\struck{(31) \quad ($\text{T}_0 3$) \quad $\forall \lambda \in Z'$, $s, t \in S$, on a}}
\]
\[
\text{\struck{$\{s, t\} \in A_0 \iff \{\lambda s, \lambda^{-1} t\} \in A_0 \implies t \neq \lambda s$}}
\]
\note{Les trois dernières lignes sont barrées de traits obliques.}

\page{110}
\note{Pagination de l'auteur : 56.}
Notons que
\[
(31) \qquad \lambda \in Z', \ s \in S \implies \{s, \lambda s\} \notin A_0 .
\]
C'est clair si $\lambda = \operatorname{id}_S$. Si on avait $\{s, \lambda s\} \in A_0$, soit $t \in S$ tel que $\{s, \lambda s, t\} \in F_0$ et soit $\rho \in G_0^{+}$ tel que $\rho(\lambda s) = \lambda s$, $\rho s = t$, on aurait donc $\underbrace{\rho(\lambda s)}_{= \lambda \rho(s)} = \lambda s$, \struck{\ill{}} $\lambda \rho(s) = \lambda s$ i.e.\ $\rho(s) = s$ donc $t = s$, absurde.

\drawing{Un triangle de sommets $s$, $\lambda s$ (en haut, reliés par un segment) et $t$ (en bas).}

On suppose maintenant
\[
(32) \quad (\text{T}_0 3) \qquad \forall \lambda \in Z', \ s, t \in S, \text{ on a } \quad
\{s, t\} \in A_0 \iff \{\lambda s, \lambda^{-1} t\} \in A_0
\]
[\emph{NB} Il suffit de prouver \struck{l'axiome} $\Longrightarrow$ car il en résulte, en l'appliquant à $(\lambda^{-1} t, \lambda s)$ au lieu de $(s, t)$, que l'on a
\[
\{\lambda^{-1} t, \lambda s\} \in A_0 \implies \lambda\{t, s\} \in A_0, \qquad \lambda\{t, s\} = \{\lambda \lambda^{-1} t, \lambda^{-1} \lambda s\} .
\]
]
\note{Sous l'accolade, l'auteur écrit « $\lambda\lambda^{-1}t$ » et « $\lambda^{-1}\lambda s$ » sous $t$ et $s$ ; le $\lambda$ devant l'accolade est surchargé. La transcription ne corrige pas l'écriture.}

D'autre part, comme $\lambda, \lambda^{-1}$ centralisent $G_0^{+}$, il suffit, pour \emph{un} repère origine $\{s_0, t_0, u_0\}$, de vérifier la condition
\[
(32\text{ bis}) \qquad \{\lambda s_0, \lambda^{-1} t_0\} \in A_0 .
\]
Si on \struck{suppose} choisit $\rho \in G_0^{+}$ tel qu'on ait

\page{111}
\note{En tête, un numéro de l'auteur biffé, illisible.}
\[
(33) \qquad t_0 = \rho s_0
\]
\drawing{Un triangle de sommets $s_0$, $t_0$ (en bas) et $u_0$ (en haut), avec un petit secteur hachuré à partir de $s_0$, figurant la rotation autour de $s_0$.}
(p.\ ex.\ $\rho$ = « rotation d'un cran autour de $u_0$, de $s_0$ vers $t_0$ »), on a donc
\[
\lambda^{-1} t_0 = \lambda^{-1} \rho s_0 = \rho \lambda^{-1} s_0 ;
\]
posons
\[
\begin{array}{l}
\lambda s_0 = g . s_0 \\
\lambda^{-1} s_0 = g^{-1} . s_0
\end{array}
\qquad \text{où } g \in \operatorname{Norm}_{G_0^{+}}(G_{0s_0}^{+}),
\]
la relation (32 bis) s'écrit aussi
\[
\{ g . s_0, \rho g^{-1} s_0 \} \in A_0
\]
ou encore
\[
\exists u \in G_0^{+}, \text{ avec } u(s_0, t_0) = (g s_0, \rho g^{-1} t_0)
\]
\note{Dans la dernière formule, l'auteur a surchargé le second terme ; la lecture « $\rho g^{-1} t_0$ » est celle de la surcharge.}
ou encore, sachant que $t_0 = \rho s_0$,
\[
u s_0 = g . s_0, \qquad u \rho . s_0 = \rho g^{-1} . s_0
\]
soit encore
\[
g^{-1} u \in G_{0s_0}^{+}, \quad g(\underbrace{\rho^{-1} u \rho}) \in G_{0s_0}^{+}
\]
donc l'axiome $\text{T}_0 3$ équivaut à :
\[
(34) \quad (\text{T}_0 3') \qquad
\begin{array}{l}
\forall g \in \operatorname{Norm}_{G_0^{+}}(G_{0s_0}^{+}), \ \exists z \in G_{0s_0}^{+} \text{ (nécessairement unique !)} \\
\text{tel qu'on ait } \ g . \operatorname{int}(\rho^{-1})(g z) \in G_{0s_0}^{+}
\end{array}
\]
\struck{On} \uncertain{Mais} \uncertain{on} \uncertain{a} \uncertain{vu} que la situation d'un 2-polyèdre orienté \struck{triangulé} \add{régulier} \uncertain{comme}, épinglé par $(s_0, t_0, u_0)$, \uncertain{équivaut} : celle des groupes $G_0^{+}$, avec des

\marginal{Dans \uncertain{l'exemple} \uncertain{qui} \uncertain{précède}, $\rho = \rho_1 = $ \uncertain{rotation} autour de l'\uncertain{origine}, \uncertain{p.\ ex.}\ \uncertain{$\{s_0, t_0\}$} \ill{} ; ou $\rho = \rho_\infty$ (rotation autour de \ill{}) \uncertain{bien} \ill{} \ill{} \ill{} \ill{} $\{s_0, t_0\}$, et}
\marginal{on a posé $g^{-1} u = z$ i.e.\ $u = g z$.}
\marginal{où $\rho$ est \struck{l'\ill{}} \uncertain{l'élément} \uncertain{de} $G_0^{+}$ fixé une fois pour toutes, tel que $\rho s_0 = t_0$ ! \uncertain{Ainsi}, $\rho = \rho_1$ ou $\rho = \rho_\infty$.}
\note{Trois notes marginales écrites en oblique le long du bord gauche ; leur lecture est en grande partie incertaine.}

\page{112}
\note{Pagination de l'auteur : 57.}
générateurs satisfaisant
\[
G_0^{+} \ni \rho_0, \rho_1, \rho_\infty \ \big| \ \rho_0^{n} = \rho_1^{2} = \rho_\infty^{3} = 1
\]
(rotations autour de $s_0$, du centre de $\{s_0, t_0\}$, et du centre de la face $\{s_0, t_0, u_0\}$), et que $G_{0s_0}^{+}$ est le sous-groupe engendré par $\rho_0$, la condition $(\text{T}_0 3')$ est donc une traduction de $\text{T}_0 3$ en termes de théorie des groupes pure. \struck{\ill{}}

\marginal{\emph{NB} \uncertain{on} \uncertain{est} \uncertain{certain} \uncertain{que} l'ordre de $\rho_0$ dans $G_0^{+}$ \ill{}}

Dans les cas qui nous intéressent, où $P_0$ est régulier \uncertain{et} \uncertain{orienté}, \uncertain{i.e.}\ \uncertain{tout} \uncertain{est} \uncertain{fixé} \uncertain{par} \struck{\ill{}} $Z'$ non orienté, on a \uncertain{connaissance} \uncertain{grâce} \uncertain{à} : $G_0$ tout entier, non seulement à $G_0^{+}$, et on pourra, dans l'interprétation précédente, choisir un $\rho \in G_0$ -- p.\ ex.\ on pourra prendre la \struck{symétrie} \add{réflexion} $\tau_0$ \struck{\ill{}} qui échange $s_0$ et $t_0$, la formule (34) devient
\[
(34') \qquad g . \operatorname{int}(\tau_0)(g z) \in \langle \rho_0 \rangle
\quad \text{\small (sous-groupe engendré par $\rho_0$)}
\]
pour $z \in \langle \rho_0 \rangle$ convenable (dépendant du choix de $g \in \operatorname{Norm}_{G_0^{+}}(\langle \rho_0 \rangle)$).]
\note{Le crochet fermant final, en bas à droite, n'a pas d'ouverture visible sur la page ; le numéro « (34) » de la ligne « (34) devient » est surchargé.}

\page{113}
Considérons maintenant, pour $\lambda \in Z'$
\[
\begin{aligned}
(35) \qquad A_\lambda &= \bigl\{ a \in \mathfrak{P}_2(S) \bigm| \exists (s, t) \in S^2, \text{ avec } \{s, t\} \in A_0, \ a = \{s, \lambda t\} \bigr\} \\
&= \bigl\{ \{s, \lambda t\} \bigm| (s, t) \in S^2, \text{ t.q. } \{s, t\} \in A_0 \bigr\}
\end{aligned}
\]
\marginal{\emph{NB} cette condition implique bien, par (31), que $s \neq \lambda t$ donc $\{s, \lambda t\} \in \mathfrak{P}_2(S)$.}
On a donc (pour $\lambda$ choisi !) \struck{$A_\lambda$} une partie de $\mathfrak{P}_2(S)$, définissant sur $S$ une structure de graphe strict, notée $\Gamma_\lambda$, qui est liée à la structure de graphe $\Gamma_0 = (S, A_0)$ associée à $P_0 = (S, A_0, F_0, \ldots)$, via $\lambda$, par le fait que $\lambda$ induit une bijection
\[
(36) \qquad
\underset{\substack{\text{étoile de} \\ s \text{ pour } \Gamma_0}}{A_0(s)}
\xrightarrow{\ \sim\ }
\underset{\substack{\text{étoile de} \\ s \text{ pour } \Gamma_\lambda}}{A_\lambda(s)}
\]
puisqu'on a, pour $(s, t) \in S^2$,
\[
(37) \qquad
\underset{\text{i.e.\ } t \in A_0(s)}{\{s, t\} \in A_0}
\iff
\underset{\text{i.e.\ } \lambda t \in A_\lambda(s)}{\{s, \lambda t\} \in A_\lambda} ,
\]
ce qui est essentiellement la définition (35), à cela près que celle-ci signifie que $\{s, \lambda t\} \in A_\lambda$ \uncertain{se} \uncertain{traduit} \uncertain{par} $\{s, t\} \in A_\lambda$ \emph{ou} $\{\lambda^{-1} s, \lambda t\} \in A_\lambda$, mais par $\text{T}_0 3$ les deux conditions sont équivalentes
\note{Sous la ligne, l'auteur glose $\{s, \lambda t\} = \{\lambda t, s\}$ et $= \{\lambda t, \lambda(\lambda^{-1} s)\}$, avec le renvoi « $\text{T}_0 3$ ». La phrase se poursuit p.~114.}

\page{114}
\note{Pagination de l'auteur : 58.}
On voit donc que pour tout $s \in S$, l'ordre du sommet $s$ pour le graphe $A_\lambda$ est \struck{\ill{}} égal à son ordre $n$ pour $A_0$. Notons aussi que par transport de structure, $G_0^{+}$ \struck{opère} transforme $A_\lambda$ en lui-même, i.e.\ $G_0^{+}$ opère par automorphismes du graphe $\Gamma_\lambda$ -- \struck{ainsi, donc} on voit aussi qu'il est \struck{transitif} \add{\uncertain{simplement}} \add{transitif} sur l'ens.\ des arêtes orientées de $\Gamma_\lambda$, puisqu'il en est ainsi pour $\Gamma_0$.

On veut exprimer que $\Gamma_\lambda$ définit une structure de polyèdre régulier triangulé \uncertain{non} \uncertain{-} \uncertain{orienté} : polyèdre isomorphe à $P_0$. On va supposer pour ceci, d'abord, que la structure de graphe $\Gamma_0 = (S_0, A_0)$ détermine la structure de polyèdre triangulé $P_0 = (S_0, F_0)$ ; i.e.\
\[
(38) \quad (\text{T}_0 4) \qquad F_0 = \bigl\{ f \in \mathfrak{P}_3(S) \bigm| \forall a \in \mathfrak{P}_2(f), \text{ on a } a \in A_0 \bigr\}
\]
(la structure de graphe détermine la structure triangulaire)
\marginal{serait-ce \uncertain{conséquence} de $\text{T}_0 1$, $\text{T}_0 2$, $\text{T}_0 3$ ?}

J'ignore, s'il existe une structure de 2-polyèdre régulier \struck{orientable} triangulé orientable, qui ne
\note{La phrase se poursuit p.~115.}

\page{115}
soit \uncertain{aussi} décrite par sa structure de graphe -- déjà -- \uncertain{d'ailleurs} \uncertain{même} sans supposer la structure orientable ni régulière, ce n'est pas évident qu'il y ait un contre-exemple. \struck{\ill{}} Mais si -- dans l'\emph{icosaèdre \uncertain{gauche}} (régulier, mais pas orientable), $S$ de cardinal $6$, \struck{\ill{}} toute $a \in \mathfrak{P}_2(S)$ est une arête (il y en a $15$), mais seulement \ill{} des $20$ éléments $f \in \mathfrak{P}_3(S)$ sont des faces…
\note{Les nombres sont lus tels qu'écrits : « cardinal 6 », « 15 » et « 20 ». Le mot qualifiant « icosaèdre », souligné, est mal formé.}

\struck{Rajoutons} On peut encore \struck{\ill{}} \add{exiger} pour \uncertain{une} structure de graphe strict, voir pour une structure de graphe $\Gamma_0 = (S_0, A_0)$, à quelles conditions l'ens.\ $F_0 \subset \mathfrak{P}_3(S)$ qu'il définit est une structure de 2-polyèdre. On trouve ceci :
(39)
\begin{itemize}
\item A1) $\forall a = \{s, t\} \in A_0$, il existe exactement deux éléments $u \in S$, tels que l'on ait $\{s, u\}, \{t, u\} \in A_0$ i.e.\ adjacents à la fois aux deux extrémités de l'arête -- i.e.\ $A_0(s) \cap A_0(t)$ est de cardinal~2.
\item A2) $\forall s \in S$, la structure de contour dont l'ens.\ des sommets est $A_0(s) = \operatorname{Et}_{\Gamma_0}(s)$ (ens.\ des sommets adjacents), et dont les arêtes sont les $a \in \mathfrak{P}_2(A_0(s))$ telles que $a \in A_0$, est \emph{connexe} non vide \struck{(\emph{NB} \ill{})}
\end{itemize}
\marginal{\emph{NB} Le fait qu'il existe un, unique tel $u$ ; \uncertain{indique} \uncertain{déjà} que $A_0$ se déduit de $F_0$ par $A_0 = \{ a \in \mathfrak{P}_2(S) \mid \exists f \in F_0, f \supset a \}$.}
\marginal{\emph{NB} Si $A_0(s) \neq \emptyset$, il résulte déjà de A1 que $\operatorname{card} A_0(s) \geq 3$.}
\note{L'en-tête « (39) » est surchargé sur un numéro antérieur ; l'indice de « $\operatorname{Et}$ » se lit $\Gamma_0$ ou $S_0$.}

\page{116}
\note{Pagination de l'auteur : 59.}
Ceci dit, on peut se poser la question si $A_\lambda$, déduit de $A_0$ via $\lambda$ comme dit dans (35), définit bien une structure de graphe satisfaisant à A1) A2), i.e.\ permettant de définir une structure 2-polyédrale (qui sera automatiquement régulière, $G_0^{+}$ formant un groupe d'automorphismes \struck{transitif} \add{\uncertain{simplement}} transitif sur l'ens.\ des arêtes orientées du dit polyèdre…).

Considérons d'abord A1, soit donc \struck{$a \in \mathfrak{P}_2(S)$} $a \in A_\lambda$, on peut écrire
\[
a = \{s, \lambda t\} \qquad \text{avec } \{s, t\} \in A_0
\]
soit $v \in S$, donc $v$ de la forme $\lambda u$, on aura
\[
\begin{aligned}
\text{\struck{$\{s, \lambda t, v\}$}}\ a \cup \{v\} \in F_\lambda
&\iff \{s, \lambda u\} \in A_\lambda, \ \{\lambda t, \lambda u\} \in A_\lambda \\
\text{i.e.\ } \{s, v\} \in A_\lambda, \ \{\lambda t, v\} \in A_\lambda \qquad
&\iff \{s, u\} \in A_0, \ \{t, \lambda u\} \in A_0
\end{aligned}
\]
\note{Les deux équivalences sont transcrites telles qu'écrites ; dans $\{s, \lambda u\}$ le $\lambda$ est surchargé, et la lecture $v$ / $u$ des lettres est incertaine par endroits.}
donc la condition A1 pour $\Gamma_\lambda$ s'écrit

\struck{Pour $(s, t) \in S^2$, tel que $\{s, t\} \in$}
(40) Pour deux sommets $s, t$ adjacents pour $P_0$, et $\lambda \in Z'$, il existe exactement deux sommets $u$ tels que \struck{l'on ait} $u$ soit adjacent à $t$, et $s$, $\lambda u$ adjacents, i.e.
\[
\operatorname{card} \lambda . A_0(s) \cap A_0(t) = 2 \quad \text{si } \{s, t\} \in A_0, \ \lambda \in Z'
\]
\note{Le numéro « (40) » est surchargé ; la reformulation verbale de (40) est confuse sur la page (« $u$ soit adjacent à $t$ » et « $s$, $\lambda u$ adjacents ») et elle est transcrite telle quelle.}

\page{117}
Cette condition n'a pas l'air tautologique ! D'autre part, la condition de connexité pour $A_\lambda(s)$ peut s'examiner à l'aide de la bijection canonique (36), induite par $\lambda$
\[
A_0(s) \xrightarrow{\ \sim\ } A_\lambda(s) ;
\]
on note que $\{\lambda t, \lambda t'\} \in A_\lambda$ ssi $\{t, \lambda t'\} \in A_0$ -- pourquoi cette relation (symétrique) entre deux éléments $t, t'$ de $A_0(s)$ définirait-elle « connexe » ? En fait, on va expliciter ces propriétés par un axiome qui, contrairement à l'axiome « par acquit de conscience » $\text{T}_0 4$, a l'air des plus substantielles :
\[
(41) \quad (\text{T}_0 5) \qquad
\begin{array}{l}
\forall \lambda \in Z', \text{ la structure de graphe } \Gamma_\lambda = (S, A_\lambda) \text{ (cf.\ (35)) est} \\
\text{\emph{isomorphe} à } \Gamma_0 = (S, A_0),
\end{array}
\]
en d'autres termes, on demande qu'il existe une permutation $g = g_\lambda \in \mathfrak{S}_S$ telle que l'on ait
\[
(42) \qquad \forall s, t \in S, \text{ on a } \quad
\{s, t\} \in A_0 \iff
\underbrace{\{g s, \lambda^{-1} g t\} \in A_0}_{\text{ce qui exprime } \{g s, g t\} \in A_\lambda}
\]
\marginal{Ici, on voit bien, sur des exemples, que l'hypothèse $n \geq 6$ \uncertain{intervient} (\ill{} \uncertain{l'oubli} \ill{} \ill{} \ill{} $n \neq 4$, $n \neq 5$ \uncertain{toujours} \ill{}) \uncertain{en} \uncertain{jeu} !}
\note{Note marginale en oblique au bord gauche, lecture en grande partie incertaine ; le $\mathfrak{S}_S$ de (42) est écrit comme un grand $\mathfrak{S}$ ou $G$ d'indice $S$.}

\page{118}
\note{Pagination de l'auteur : 60.}
Du coup il est clair que $A_\lambda$ définit une structure de 2-polyèdre \emph{régulier} et tout
\[
(43) \qquad P_\lambda = (S, \underset{\cap\ \mathfrak{P}_2(S)}{A_\lambda}, \underset{\cap\ \mathfrak{P}_3(S)}{F_\lambda}, \text{relations d'incidence évidentes}) .
\]
En fait, on a vu que $G_0^{+}$ opère par automorphismes de $\Gamma_\lambda$, donc il opère par automorphismes de $P_\lambda$ -- on devine que ce sont les automorphismes directs -- mais on voudrait que $G_0$ aussi opère par automorphismes, pour ceci on va supposer
\[
(44) \quad (\text{T}_0 6) \qquad
\begin{array}{l}
Z' \text{ commute (non seulement à } G_0^{+}, \text{ mais) à tout le groupe } G_0,
\end{array}
\]
ce qui implique que $G_0$ opère sur $S$ par $\Gamma_\lambda$-automorphismes, donc par $P_\lambda$-automorphismes. \struck{On} J'aimerais vérifier que l'homomorphisme canonique injectif
\[
G_0 \hookrightarrow \operatorname{Aut}(P_\lambda) \overset{\text{déf}}{=} G_\lambda
\]
est un isomorphisme -- comme $P_\lambda \simeq P_0$, on a $G_\lambda \simeq G_0$, et l'isomorphisme est évident si
\marginal{On pourrait aussi exiger que la condition (42) de $\text{T}_0 5$ soit vérifiée par un \struck{\ill{}} $g \in G_0$ \uncertain{respectant} \uncertain{la} \uncertain{structure} de $P_0$, \uncertain{plutôt} \uncertain{qu'un} \ill{} \uncertain{arbitraire} de \ill{} \ill{} $\mathfrak{S}_S$ \ill{}, i.e.\ \ill{} \ill{} \ill{} $Z'$ \ill{} \ill{} \ill{} \ill{} $\text{T}_0 2$.}
\note{L'en-tête « (44) » est surchargé sur un « (42) » ; la longue note marginale oblique de gauche n'est lue qu'en partie, et sa lecture est très incertaine. La phrase finale se poursuit p.~119.}

\page{119}
$G_0$ est fini i.e.\ $S$ fini (sinon on n'est pas tellement clair sur les définitions). On peut aussi \uncertain{déduire} que \struck{pour} $G_0^{+}$ est simplement transitif sur l'ens.\ $\vec A_\lambda$ des arêtes orientées de \struck{$A$} $P_\lambda$, donc \emph{si \uncertain{on} \uncertain{savait}} \uncertain{que} $G_0$ est transitif (donc \struck{simplement transitif}) sur \struck{l'ens.\ des repères de $P_\lambda$,} \uncertain{l'un} des repères de $A_\lambda$, \struck{$r_0 = (s_0, t_0, u_0)$} \struck{$r_1 = (s_1, t_1, u_1)$} \struck{il} un \struck{\ill{}} il suffit de prouver que pour un \uncertain{repère} $(s_0, t_0, u_0)$, (donc pour tout repère), il existe un $u \in G_0$ \struck{\ill{}} qui fixe $s_0$, $t_0$, et non $u_0$, \struck{\uncertain{Mais} \uncertain{ceci}} \uncertain{n'est} \uncertain{pas} l'identité. Mais s'il en était ainsi, alors signifierait que $G_0$ est simplement transitif sur l'ens.\ $\vec A_\lambda$ (le stabilisateur d'un $\vec a = (s_0, t_0) \in \vec A_\lambda$ dans $G$ se réduisant à l'élément neutre),
ce qui contredit le fait que $G_0^{+}$ est simplement transitif et \struck{\ill{}} $G_0^{+} \subsetneqq G_0$. Reste à bien voir
\[
(43) \qquad G_0 \xrightarrow{\ \sim\ } \operatorname{Aut}(P_\lambda) = G_\lambda ,
\]
on aimerait prouver que cet iso.\ induit
\[
(44) \qquad G_0^{+} \xrightarrow{\ \sim\ } G_\lambda^{+} \qquad ?
\]
\note{Les numéros « (43) » et « (44) » répètent ceux des p.~118 ; ils sont transcrits tels quels.}

\page{120}
\note{Pagination de l'auteur : 61.}
Cela se voit facilement si l'ordre des points de $P_0$ est impair, car alors $G_0^{+}$ se décrit intrinsèquement comme \ill{} au \uncertain{terme} de $G_0$, par le fait que c'est le \emph{seul} sous-groupe d'indice~2. Si par contre $n$ est pair, ce n'est pas évident -- il y a \emph{trois} a priori sous-groupes d'indice~2 possibles, car $(G_{\infty 3})_{\mathrm{ab}} \simeq \mathbb{F}_2 \times \mathbb{F}_2$. Mais il n'y en a qu'un seul qui contienne $\rho_1$, donc la question est si $\rho_1^{r}$ (relatif à un repère $r = \{s_0, t_0, u_0\}$ de $P_0$) \uncertain{respecte} l'orientation \struck{dans} \uncertain{de} $P_\lambda$…
\note{L'indice du groupe dans $(G_{\infty 3})_{\mathrm{ab}}$ est lu tel qu'il se présente (« $\infty 3$ », puis « ab ») ; lecture incertaine.}
\marginal{il faudra y revenir ailleurs…}

Je m'intéresse maintenant aux structures $P_\lambda$, $\lambda \in Z'$, \struck{Notons d'abord} les structures polyédrales triangulées satisfaisant aux axiomes $(\text{T}_0 1)$ à $(\text{T}_0 6)$ seront appelées (à défaut de mieux, et provisoirement) « \emph{spéciales} ». Je dirais que deux telles structures, \add{définies par des ens.\ d'arêtes} $A_0, A_1 \subset \mathfrak{P}_2(S)$, sur un même ensemble $S$ sont \emph{équivalentes}, s'il existe un $\lambda \in \mathfrak{S}_S$, commutant à $\operatorname{Aut}(S, F_0)$, de telle façon que l'on ait
\note{La phrase et l'argument se poursuivent au-delà de la p.~120, hors de ce lot.}

\end{document}
